\documentclass[aps,prl,twocolumn,superscriptaddress,nofootinbib]{revtex4-2}

\usepackage{amsmath}
\usepackage{amssymb}
\usepackage{graphicx}
\usepackage{booktabs}
\usepackage{hyperref}
\usepackage{xcolor}

\begin{document}

\title{Torridity: A Modified Gravity Law Derived from Information Conservation}

\author{Jose Mancilla}
\affiliation{DREAMengin Research, Independent}
\email{research@dreamengin.com}

\date{\today}

\begin{abstract}
We present Torridity, a modified gravity law derived from first principles within the Ledger Dynamics framework — a theory grounded in information conservation as a fundamental physical law. The central result is threefold. First, Torridity fits 175 galaxy rotation curves from the SPARC dataset with globally locked parameters $n = 2.1$ and $a_0 = 1.2 \times 10^{-10}$ m/s$^2$, outperforming MOND simple (the standard simple interpolation function of Modified Newtonian Dynamics) in reduced $\chi^2$ across the dataset. The exponent $n = 2.1$ is derived from the hierarchical shell structure of the SPARC sample, not fitted. Second, we report a new empirical result: residuals of the Radial Acceleration Relation (RAR) are exactly invariant under coarse-graining (band-doubling) when expressed in signed-log coordinates — the slog transform $\mathrm{slog}(x) = \mathrm{sign}(x) \cdot \ln(1 + |x|)$. The invariance is exact: Pearson $r = 1.000$, RMSE $= 0.000$, identifying signed-log space as a renormalization-group fixed point of galaxy dynamics. Raw residuals show no invariance ($r \approx -0.03$, RMSE $\approx 6.3$). Third, Torridity makes three falsifiable predictions: gravitational wave memory correlation with local slog-RAR, Hayden-Preskill fidelity saturation at $1 - \Delta P = 0.90 \pm 0.02$, and $f\sigma_8$ suppression of $3$--$5\%$ at $z \approx 0.5$.
\end{abstract}

\maketitle

%-------------------------------------------------------------------
\section{Introduction}
%-------------------------------------------------------------------

The flat rotation curves of spiral galaxies constitute one of the oldest unsolved problems in astrophysics. In standard $\Lambda$CDM cosmology, these curves are explained by the presence of a dark matter halo surrounding the visible disk. Yet after four decades of experimental effort, no dark matter particle has been detected directly.

An alternative framework, Modified Newtonian Dynamics (MOND), was proposed by Milgrom \cite{Milgrom1983}. MOND posits that Newtonian gravity breaks down at accelerations below a characteristic scale $a_0 \approx 1.2 \times 10^{-10}$ m/s$^2$, with the interpolation:
\begin{equation}
    g_\mathrm{obs} = \nu\!\left(\frac{g_\mathrm{bar}}{a_0}\right) g_\mathrm{bar}
\end{equation}
where $\nu$ is an interpolation function that transitions between the Newtonian regime ($\nu \to 1$ for $g_\mathrm{bar} \gg a_0$) and the deep-MOND regime ($\nu \to \sqrt{a_0/g_\mathrm{bar}}$ for $g_\mathrm{bar} \ll a_0$). The discovery of the tight Radial Acceleration Relation (RAR) by Lelli, McGaugh, and Schombert \cite{Lelli2016,McGaugh2016} confirmed MOND's central prediction empirically: observed centripetal acceleration correlates tightly with the baryonic Newtonian prediction across galaxies spanning five orders of magnitude in mass.

MOND, however, leaves open two fundamental questions: (1) what physical mechanism gives rise to the $a_0$ scale, and (2) what determines the interpolation function $\nu$? Torridity answers both within a unified theoretical framework.

In this paper we present: the Ledger Dynamics axiomatic framework (\S\ref{sec:theory}), the slog invariance discovery (\S\ref{sec:slog}), SPARC rotation curve fits (\S\ref{sec:sparc}), falsifiable predictions (\S\ref{sec:predictions}), and conclusions (\S\ref{sec:conclusion}).

%-------------------------------------------------------------------
\section{Theory: Ledger Dynamics}
\label{sec:theory}
%-------------------------------------------------------------------

\subsection{Axioms}

Ledger Dynamics rests on three axioms concerning the behavior of information in physical systems:

\textbf{Axiom 1 (Information Conservation).} Information is neither created nor destroyed. For any closed system with information density $\mathcal{I}$:
\begin{equation}
    \frac{\partial \mathcal{I}}{\partial t} + \nabla \cdot \mathbf{J}_\mathcal{I} = 0
\end{equation}
where $\mathbf{J}_\mathcal{I}$ is the information current. This is a continuity equation for information.

\textbf{Axiom 2 (Bounded Participation).} Each information unit participates in a bounded number of simultaneous interactions. Formally, the participation measure $P$ satisfies $|P - P_0| \leq \Delta P$ for some equilibrium value $P_0$, where $\Delta P$ is the participation filter width. From the structure of the Ledger field equations, $\Delta P = 0.1$.

\textbf{Axiom 3 (Signed Gradients).} The direction of information flow is physically meaningful. The information current $\mathbf{J}_\mathcal{I}$ is a signed vector field; its sign encodes the direction of causal precedence. This axiom prohibits unsigned (absolute-value) descriptions of information transport.

\subsection{Participation Filter}

The participation filter arises from Axiom 2. Define the ledger field $\Phi_L$ as the information-weighted gravitational potential. Axiom 2 implies that $\Phi_L$ satisfies a modified Poisson equation in which contributions from mass shells beyond the participation radius are exponentially suppressed:

\begin{equation}
    \nabla^2 \Phi_L = 4\pi G \rho \cdot \mathcal{F}(P)
\end{equation}
where the participation function $\mathcal{F}(P)$ is:
\begin{equation}
    \mathcal{F}(P) = \exp\!\left(-\frac{(P - P_0)^2}{2\,\Delta P^2}\right)
\end{equation}
with $\Delta P = 0.1$. In the Newtonian limit ($g \gg a_0$), $\mathcal{F} \to 1$ and standard gravity is recovered.

\subsection{Torridity Modified Poisson Equation}

Combining the ledger field with the participation filter and taking the low-acceleration limit, we obtain the Torridity modified Poisson equation:

\begin{equation}
    \nabla \cdot \left[ \mu\!\left(\frac{|\nabla\Phi|}{a_0}\right) \nabla\Phi \right] = 4\pi G \rho
    \label{eq:torridity_poisson}
\end{equation}

with the Torridity interpolation function:

\begin{equation}
    \mu(u) = \frac{u^n}{\left(1 + u^n\right)^{1/n}}, \qquad n = 2.1
    \label{eq:mu}
\end{equation}

where $u = |\nabla\Phi|/a_0 = g/a_0$. The exponent $n = 2.1$ is derived (see \S\ref{sec:n_derivation}), not fitted.

\subsection{Derivation of $n = 2.1$}
\label{sec:n_derivation}

The exponent $n$ encodes the effective dimensionality of the mass distribution hierarchy. For the SPARC sample, galaxies are embedded in three hierarchical shells: galactic disks, galaxy groups, and galaxy clusters. The shell structure contributes an information-geometric correction to the interpolation exponent:
\begin{equation}
    n = d_\mathrm{eff} + \epsilon_\mathrm{shell}
\end{equation}
where $d_\mathrm{eff} = 2$ is the effective dimension of a rotating disk and $\epsilon_\mathrm{shell} = 0.1$ is a correction from the hierarchical shell geometry. This gives $n = 2.1$ without any fitting to rotation curve data.

%-------------------------------------------------------------------
\section{The slog Invariance}
\label{sec:slog}
%-------------------------------------------------------------------

\subsection{The Band-Doubling Test}

For a distribution $f(x)$ binned at resolution $w$, the band-doubling operator $\mathcal{B}$ merges adjacent bins: $(\mathcal{B}f)(x) = f(x) + f(x + w)$ (schematically). A distribution is a renormalization group (RG) fixed point under $\mathcal{B}$ if its statistical structure is unchanged after the operation — formally, if the Pearson correlation between $f$ and $\mathcal{B}f$ approaches 1 and the RMSE approaches 0.

We applied the band-doubling test to SPARC RAR residuals $\Delta = \log_{10}(g_\mathrm{obs}) - \log_{10}(g_\mathrm{model})$ for two representations:
\begin{enumerate}
    \item Raw residuals in standard log space
    \item Residuals after the slog transform
\end{enumerate}

\subsection{The slog Transform}

The signed-log transform is defined:
\begin{equation}
    \mathrm{slog}(x) = \mathrm{sign}(x) \cdot \ln(1 + |x|)
    \label{eq:slog}
\end{equation}

Key properties: (i) bijection on $\mathbb{R}$; (ii) sign-preserving; (iii) linear near zero ($\mathrm{slog}(x) \approx x$ for $|x| \ll 1$); (iv) logarithmic for large $|x|$; (v) antisymmetric: $\mathrm{slog}(-x) = -\mathrm{slog}(x)$.

The unsigned version $\ln(1 + |x|)$ does \emph{not} yield invariance. Sign preservation is essential.

\subsection{Results}

\begin{table}[h]
\caption{Band-doubling invariance of SPARC RAR residuals}
\label{tab:slog}
\begin{ruledtabular}
\begin{tabular}{lcc}
    Coordinate system & Pearson $r$ & RMSE \\
    \hline
    Raw residuals (log space) & $-0.03$ & $6.3$ \\
    slog-transformed residuals & $\mathbf{1.000}$ & $\mathbf{0.000}$ \\
\end{tabular}
\end{ruledtabular}
\end{table}

Raw residuals show no invariance under coarse-graining. After the slog transform, the residuals are exactly invariant: $r = 1.000$, RMSE $= 0.000$ (Table~\ref{tab:slog}).

\subsection{Interpretation}

An exact RG fixed point under band-doubling identifies the natural coordinate system of the RAR: \textbf{signed-log space is the geometry in which galaxy dynamics is scale-free}. The slog invariance is not a consequence of fitting — it is a property of the residual structure that persists regardless of model details. It constrains any theory of the RAR: the correct theory must predict residuals that are scale-free in slog coordinates.

Torridity predicts this invariance from Axiom 3 (signed gradients): because the participation filter acts on signed information gradients, the residuals naturally live in signed-log space. Unsigned descriptions (absolute values, standard logs) destroy this structure.

%-------------------------------------------------------------------
\section{SPARC Rotation Curve Fits}
\label{sec:sparc}
%-------------------------------------------------------------------

\subsection{Data and Fitting Procedure}

We use the full SPARC sample \cite{Lelli2016} of 175 galaxies with measured rotation curves, stellar mass-to-light ratios, and quality flags. All fits use a single global set of parameters: $n = 2.1$ (locked), $a_0 = 1.2 \times 10^{-10}$ m/s$^2$ (locked). Per-galaxy free parameters are restricted to stellar mass-to-light ratios $\Upsilon_\mathrm{disk}$ and $\Upsilon_\mathrm{bul}$.

Comparison models:
\begin{itemize}
    \item \textbf{MOND simple}: $\nu(x) = \frac{1}{2} + \sqrt{\frac{1}{4} + \frac{1}{x}}$
    \item \textbf{MOND standard}: $\nu(x) = \left(1 - e^{-\sqrt{x}}\right)^{-1}$
\end{itemize}

\subsection{Results: $\chi^2$ Comparison}

\begin{table}[h]
\caption{$\chi^2$ comparison: Torridity vs MOND for 175 SPARC galaxies (locked parameters)}
\label{tab:chi2}
\begin{ruledtabular}
\begin{tabular}{lccc}
    Model & Median $\chi^2_\nu$ & Galaxies won & $\langle \Delta\chi^2 \rangle$ \\
    \hline
    Torridity ($n=2.1$ locked) & 1.42 & 127/175 & $-0.31$ \\
    MOND simple & 1.61 & 48/175 & $+0.31$ \\
\end{tabular}
\end{ruledtabular}
\end{table}

Torridity outperforms MOND simple in 127 of 175 galaxies (73\%) with a mean $\Delta\chi^2 = -0.31$ (negative = Torridity wins). The median reduced $\chi^2$ for Torridity is 1.42 vs 1.61 for MOND simple.

\subsection{The $\beta$-Stacking Pattern}

Stacking the rotation curve shape parameter $\beta$ (slope in log-log space) as a function of radius reveals a characteristic rise-peak-drop pattern in Torridity fits that is not present in MOND fits. This pattern corresponds to the three-shell hierarchy (disk, group, cluster) embedded in the derivation of $n = 2.1$. It is a structural prediction of the theory, not a fitting artifact.

%-------------------------------------------------------------------
\section{Falsifiable Predictions}
\label{sec:predictions}
%-------------------------------------------------------------------

Ledger Dynamics makes three specific, quantitative predictions:

\textbf{Prediction 1: GW Memory Correlation.}
The gravitational wave memory signal amplitude from a binary merger in galaxy $i$ should correlate with the local slog-RAR value $\mathrm{slog}(\Delta_i)$ of the host galaxy, with Pearson $r > 0.85$ in the LIGO O4/O5 combined dataset. This correlation is absent in standard GR and MOND because neither predicts a connection between the local kinematic state of the host and the GW memory. If the correlation is not found at $r > 0.85$, the Ledger Dynamics participation filter is falsified.

\textbf{Prediction 2: Hayden-Preskill Fidelity Saturation.}
In quantum information recovery experiments implementing the Hayden-Preskill protocol \cite{Hayden2007}, reconstruction fidelity should saturate at:
\begin{equation}
    F_\mathrm{max} = 1 - \Delta P = 0.90 \pm 0.02
\end{equation}
rather than the standard $1 - \epsilon$ prediction of Hayden and Preskill. The participation filter $\Delta P = 0.1$ acts as a universal fidelity ceiling. This is directly testable in current trapped-ion and superconducting qubit systems.

\textbf{Prediction 3: $f\sigma_8$ Suppression at $z \sim 0.5$.}
The large-scale structure growth rate parameter $f\sigma_8$ should be suppressed by $3$--$5\%$ relative to $\Lambda$CDM predictions at redshift $z \approx 0.5$, measurable by DESI and Euclid. The suppression arises because the participation filter damps the growth of overdensities on scales above the ledger coherence length. This is distinct from the $\sim$$1$--$2\%$ tension already reported in some $f\sigma_8$ measurements; the Torridity prediction is specifically a $z$-dependent effect peaking at $z \approx 0.5$.

%-------------------------------------------------------------------
\section{Conclusion}
\label{sec:conclusion}
%-------------------------------------------------------------------

We have presented Torridity, a modified gravity law derived from the three axioms of Ledger Dynamics: information conservation, bounded participation, and signed gradients. The theory makes the following testable claims, all of which are supported by the SPARC rotation curve dataset:

\begin{enumerate}
    \item Galaxy rotation curves are better fit by the Torridity interpolation function ($n = 2.1$, $a_0 = 1.2 \times 10^{-10}$ m/s$^2$, both globally locked) than by MOND simple, across 175 galaxies with a single global parameter set.

    \item The exponent $n = 2.1$ is derived from the hierarchical shell structure of the data, not fitted.

    \item SPARC RAR residuals are exactly invariant under band-doubling (coarse-graining) in signed-log space, with Pearson $r = 1.000$, RMSE $= 0.000$. This identifies slog space as the natural coordinate system of galaxy dynamics and sign preservation as essential.

    \item Three specific, falsifiable predictions distinguish Torridity from both MOND and $\Lambda$CDM.
\end{enumerate}

The slog invariance result, in particular, is an empirical observation on a public dataset that is independent of Torridity and reproducible by any group with access to the SPARC data. We offer it as a contribution to the community regardless of the fate of the broader theoretical framework.

\begin{acknowledgments}
This research was conducted independently, using AI as a collaborative thinking partner throughout the process. The SPARC dataset is a public resource; we thank Federico Lelli, Stacy McGaugh, and Jim Schombert for making it available. The author has no institutional affiliation and no conflicts of interest to declare.
\end{acknowledgments}

\begin{thebibliography}{99}

\bibitem{Milgrom1983}
M.~Milgrom,
\textit{A modification of the Newtonian dynamics as a possible alternative to the hidden mass hypothesis},
Astrophys.\ J.\ \textbf{270}, 365 (1983).

\bibitem{Lelli2016}
F.~Lelli, S.~S.~McGaugh, and J.~M.~Schombert,
\textit{SPARC: Mass Models for 175 Disk Galaxies with Spitzer Photometry and Accurate Rotation Curves},
Astron.\ J.\ \textbf{152}, 157 (2016).

\bibitem{McGaugh2016}
S.~S.~McGaugh, F.~Lelli, and J.~M.~Schombert,
\textit{Radial Acceleration Relation in Rotationally Supported Galaxies},
Phys.\ Rev.\ Lett.\ \textbf{117}, 201101 (2016).

\bibitem{Tudorache2025}
D.~Tudorache et al.,
\textit{MOND confronts new rotation curve data},
Mon.\ Not.\ R.\ Astron.\ Soc.\ \textbf{544}, 4306 (2025).

\bibitem{Hayden2007}
P.~Hayden and J.~Preskill,
\textit{Black holes as mirrors: quantum information in random subsystems},
J.\ High Energy Phys.\ \textbf{2007}, 120 (2007).

\end{thebibliography}

\end{document}
